\section{FishBot development background}
\label{sec:background}

\sestina{} is the seventh configuration in a lineage, and several of its components and most of its
evaluation machinery exist because of what earlier versions found. This section is deliberately
short. Its purpose is to explain why \sestina{} looks the way it does, not to re-report earlier
studies; the detailed corrections those studies require are collected in Appendix~\ref{app:corrections}.

\subsection{Engine, inference, and policy foundations (v0.2--v0.4)}
\label{sec:background-early}

The first versions established the domain and the inference. A count-conditioned, empirically tuned
linear policy in a scripting-language simulator became a C++ engine implementing the specified
rules and carrying an exact constraint-conditioned reference distribution, a faster deployed
approximation, a twenty-feature ask score and a value-based declaration rule. The durable
contribution of that period is the inference object of \S\ref{sec:inference}: an exactly computable
observer-conditioned constraint distribution over deals, cheap enough to run in all six seats at
once.

Earlier evaluations used smaller samples, unpaired ablations, and point-estimate comparisons: cells
were a few thousand games against a handful of opponents.

\subsection{Deadlock mitigation (v0.5)}
\label{sec:background-vfive}

FishBot \vfive{} addressed a mirror-play deadlock in which two agents repeatedly requested
unavailable cards until a forced declaration ended the sequence. Removing this behaviour did not
produce a statistically resolved win-rate improvement in the reported evaluation. The result
motivated separate checks of behavioural validity and playing strength, as well as reassessment of
the evaluation's precision.

\subsection{Policy and evaluation revisions (v0.6)}
\label{sec:background-vsix}

The v0.6 cycle set out to recover several points of headroom that the accumulated design register
described. It found little of it, and found in the process that much of the measurement machinery
lacked the precision its conclusions required. Four results from that cycle bear directly on
\sestina{}.

\paragraph{Tied action scores.} Measured in the \vfive{} mirror, the policy that cycle set out to diagnose, \vsixTieShareFive\% of
\emph{contested} ask decisions end with two or more candidates scoring bit-for-bit identically, and
the winner is whichever the enumerator emitted first. \vsixTieSameSuitFive\% of those ties are two
cards of one half-suit at one target, and the exact posterior \emph{separates} them in
\vsixTieExactSepPctFive\% of cases, meaning it assigns them equal probability in every case. The
same measurement on the \vsix{} mirror gives \vsevenTieContested\%, and that is the figure
\S\ref{sec:agent-tiebreak} and \S\ref{sec:development-tie} use, because \sestina{} modifies
\vsix{} rather than \vfive{}. Every rule that ranks tied candidates by that probability therefore
realises the same hit rate, and so does enumeration order. The largest item on the inherited design
register was not a defect but a property of the information set. What remains addressable is the
dependence on enumeration order itself, which is a correctness problem, and
\S\ref{sec:agent-tiebreak} is how \sestina{} addresses it.

\paragraph{Belief-model comparison.} Scored as predictors of card location on identical states in
the \vfive{} mirror, the ranking is:

\begin{center}\tabfont
\begin{tabular}{lrr}
\toprule
Posterior & Log loss (nats) & Argmax hit rate \\
\midrule
Exact, uniform deal prior & \vsixBeliefExactNLL{} & \vsixBeliefExactHit\% \\
Approximation, certificate half of the prior deleted ($\theta = 0$) & \vsixBeliefNoPriorNLL{} & \vsixBeliefNoPriorHit\% \\
\textbf{\vfive{}'s deployed approximation} & \textbf{\vsixBeliefShippedNLL{}} & \textbf{\vsixBeliefShippedHit\%} \\
\bottomrule
\end{tabular}
\end{center}

The exact constraint-based model has the highest log loss among the compared models, and it
remains highest with the behavioural weighting deleted outright as well ($\theta = \varphi = 0$:
\vsevenBeliefPriorlessNLL{} nats and \vsevenBeliefPriorlessHit\%), so most of the log-loss gap has
no behavioural content, though most of the hit-rate gap does. The certificate half of the weighting
is the gap between the second and third rows; the weighting as a whole is the gap from the
$\theta = \varphi = 0$ figures to the third. In the implementation, the block construction of
\S\ref{sec:inference} has no parameter through which an opponent's tendencies could enter; the
behavioural weighting enters only through the deployed paths, as the initialiser of the Sinkhorn fit
(\S\ref{sec:inference-prior}) and, in the search, as an importance weight on exact samples
(\S\ref{sec:agent-search}). Whether an exact constraint computation could be combined with a
behavioural model in one object is not addressed here. This closed an
experiment three studies had listed as outstanding, and it established a methodological rule the
project has applied since, namely that exact constraint-based inference is an assumption to be
tested rather than a property to be assumed. \S\ref{sec:development-alloc} reports the second
comparison of that kind.

\paragraph{Search-selection bias.} v0.6
built a search that samples full deals from the public posterior and plays each forward with all six
players reconstructed at their own information sets, rather than made clairvoyant, which
\S\ref{sec:game-why} explains would leave the rollout unable to discriminate between actions. Chosen by an unguarded argmax over sampled means,
that search is \vsixCurseGap{}~pp \emph{worse} than the blueprint it searches from, measured against a
control that forces the same code to return the blueprint's own choice at identical budget, seed and
sample size. The cause is the statistic rather than the search (\S\ref{sec:agent-search}), and a
paired lower-confidence-bound deviation rule removed the measured deficit in that experiment. The guarded search, which v0.6 measured and then deployed switched off, is a component
of \sestina{} (\S\ref{sec:agent-search}).

\paragraph{Evaluation resolution and optimiser design.} Pooling every same-pair cell in the v0.6 battery gave \vsixHeadPooled\% with a naive $z$ of
\vsixHeadZ{}, a null that an earlier draft of that report published as one. Re-run at high power
across \vsixBigFiveBanks{} disjoint banks, the same comparison is \vsixBigFive\% over
\vsixBigFiveN{} games $[\vsixBigFiveLo, \vsixBigFiveHi]$, above parity on all \vsixBigFiveAbove{}
banks. The effect had always been present; the instrument could not see it. In the same cycle, five
mechanisms cleared a 95\% paired interval at the per-cell budget the project had used since v0.3 and
returned exactly zero at three times that budget. Separately, the fitting harness was found to be
incapable of moving a policy at all, using one step size for thirty-four coordinates spanning two and
a half orders of magnitude, an unpaired objective, and an objective documented as a worst case that
was in fact a weighted mean.

\subsection{Development objectives}
\label{sec:background-handoff}

The v0.6 cycle left a policy roughly nine tenths of a point stronger than its predecessor, a design
register most of whose entries had not produced a resolved gain, a repaired optimiser, and the open
question of whether, at the budget the project had used, the absence of resolved gains reflected
the policy class or the precision of the instrument.

The v0.6 results motivated five development stages: evaluation calibration, adversary generation,
candidate-mechanism testing, configuration selection and freezing, and holdout evaluation under a
preregistered protocol. \S\ref{sec:development} and \S\ref{sec:results} report the exploratory and
holdout stages, respectively. The protocol separated configuration selection from final evaluation
to reduce selection-related bias; the software processes and information restrictions that separated
the stages are stated in \S\ref{sec:evaluation-seal}.
